% Transcripción LaTeX reconstruida a partir de certamen-2.pdf.
% La fuente original del PDF no se encontró entre los archivos disponibles.
\documentclass[11pt,a4paper]{article}
\usepackage[utf8]{inputenc}
\usepackage[T1]{fontenc}
\usepackage[spanish,es-nodecimaldot]{babel}
\usepackage{amsmath,amssymb,amsthm,bm}
\usepackage{geometry}
\usepackage{lmodern}
\geometry{margin=2.5cm}

\newcommand{\R}{\mathbb{R}}
\newcommand{\C}{\mathbb{C}}
\newcommand{\M}{\mathrm{M}}
\newcommand{\SU}{\mathrm{SU}}
\newcommand{\SO}{\mathrm{SO}}
\newcommand{\tr}{\operatorname{tr}}
\newcommand{\im}{\operatorname{Im}}
\newcommand{\kerp}{\operatorname{ker}}
\newcommand{\vect}[1]{\boldsymbol{#1}}
\newcommand{\set}[1]{\left\{#1\right\}}

\title{Certamen 2 --- Introducción a la Teoría de Grupos}
\author{José Ignacio Rosas Sepúlveda}
\date{10 de julio de 2026}

\begin{document}
\maketitle

\section*{Problema}

Muestre que
\[
\SU(2)/\mathbb{Z}_2\cong\SO(3).
\]

\section*{Solución}

La idea de la demostración consiste en construir un homomorfismo sobreyectivo
\[
\Phi:\SU(2)\longrightarrow\SO(3)
\]
cuya pérdida de información sea exactamente el subgrupo de dos elementos
\[
\mathbb{Z}_2=\set{I_2,-I_2}.
\]
Luego, usando el teorema fundamental de los homomorfismos, se obtendrá
\[
\SU(2)/\ker\Phi\cong\SO(3),
\]
y como $\ker\Phi=\mathbb{Z}_2$, seguirá el resultado pedido.

\subsection*{Grupos involucrados}

Recordemos que
\[
\SU(2)=\set{U\in\M_2(\C):U^\dagger U=I_2,\ \det U=1},
\]
y
\[
\SO(3)=\set{R\in\M_3(\R):R^{\mathsf T}R=I_3,\ \det R=1}.
\]
El grupo $\SU(2)$ está formado por matrices unitarias complejas de determinante uno, mientras que $\SO(3)$ está formado por rotaciones reales tridimensionales.

Dentro de $\SU(2)$ consideramos el subconjunto
\[
\mathbb{Z}_2=\set{I_2,-I_2}.
\]
Este conjunto es un subgrupo de $\SU(2)$, pues
\[
I_2I_2=I_2,\qquad I_2(-I_2)=-I_2,\qquad (-I_2)^2=I_2.
\]
Además,
\[
\det(I_2)=1,\qquad \det(-I_2)=(-1)^2\det(I_2)=1,
\]
por lo que $I_2,-I_2\in\SU(2)$. Este subgrupo es isomorfo al grupo cíclico de dos elementos. Por abuso usual de notación, escribiremos el mismo símbolo $\mathbb{Z}_2$ para este subgrupo central de $\SU(2)$. Como $I_2$ y $-I_2$ conmutan con toda matriz $U\in\SU(2)$, se tiene
\[
U\mathbb{Z}_2U^{-1}=\mathbb{Z}_2,\qquad \forall U\in\SU(2).
\]
Por lo tanto,
\[
\mathbb{Z}_2\trianglelefteq\SU(2),
\]
y el grupo cociente $\SU(2)/\mathbb{Z}_2$ está bien definido.

\subsection*{Matrices de Pauli y representación de vectores de $\R^3$}

Consideremos las matrices de Pauli
\[
\sigma_1=\begin{pmatrix}0&1\\1&0\end{pmatrix},\qquad
\sigma_2=\begin{pmatrix}0&-i\\i&0\end{pmatrix},\qquad
\sigma_3=\begin{pmatrix}1&0\\0&-1\end{pmatrix}.
\]
Estas matrices satisfacen
\[
\sigma_i\sigma_j=\delta_{ij}I_2+i\varepsilon_{ijk}\sigma_k,
\]
donde $\varepsilon_{ijk}$ es el símbolo de Levi--Civita. En particular, para $\vect a,\vect b\in\R^3$,
\[
(\vect a\cdot\vect\sigma)(\vect b\cdot\vect\sigma)
=(\vect a\cdot\vect b)I_2+i(\vect a\times\vect b)\cdot\vect\sigma.
\]

A cada vector $\vect x=(x_1,x_2,x_3)\in\R^3$ le asociamos la matriz
\[
X(\vect x)=\vect x\cdot\vect\sigma=x_1\sigma_1+x_2\sigma_2+x_3\sigma_3.
\]
En forma explícita,
\[
X(\vect x)=
\begin{pmatrix}
x_3&x_1-ix_2\\
x_1+ix_2&-x_3
\end{pmatrix}.
\]
Esta matriz es hermítica y de traza cero:
\[
X(\vect x)^\dagger=X(\vect x),\qquad \tr X(\vect x)=0.
\]
Además,
\[
X(\vect x)^2=(\vect x\cdot\vect x)I_2=|\vect x|^2I_2,
\]
y, por lo tanto,
\[
\det X(\vect x)=-|\vect x|^2.
\]
Así, la norma euclidiana de $\vect x$ queda codificada en el determinante de la matriz $X(\vect x)$.

\subsection*{Construcción del homomorfismo $\Phi:\SU(2)\to\SO(3)$}

Sea $U\in\SU(2)$. Definimos una transformación sobre las matrices $X(\vect x)$ por conjugación:
\[
X(\vect x)\longmapsto UX(\vect x)U^\dagger.
\]
Como $U$ es unitaria, se tiene $U^\dagger=U^{-1}$. Además, si $X(\vect x)$ es hermítica y de traza cero, entonces $UX(\vect x)U^\dagger$ también es hermítica y de traza cero. En efecto,
\[
\bigl(UX(\vect x)U^\dagger\bigr)^\dagger
=UX(\vect x)^\dagger U^\dagger
=UX(\vect x)U^\dagger,
\]
y
\[
\tr\bigl(UX(\vect x)U^\dagger\bigr)
=\tr\bigl(U^\dagger U X(\vect x)\bigr)
=\tr X(\vect x)=0.
\]
Por lo tanto, existe un único vector $\vect x'\in\R^3$ tal que
\[
UX(\vect x)U^\dagger=X(\vect x').
\]
Definimos entonces $\Phi(U)$ mediante la relación
\[
X\bigl(\Phi(U)\vect x\bigr)=UX(\vect x)U^\dagger,\qquad \vect x\in\R^3.
\]
La transformación $\Phi(U):\R^3\to\R^3$ es lineal, pues la conjugación por $U$ es lineal en $X(\vect x)$.

Mostremos ahora que $\Phi(U)$ preserva la norma. Usando la invariancia del determinante bajo conjugación,
\begin{align*}
-|\Phi(U)\vect x|^2
&=\det X\bigl(\Phi(U)\vect x\bigr)\\
&=\det\bigl(UX(\vect x)U^\dagger\bigr)\\
&=\det(U)\det X(\vect x)\det(U^\dagger)\\
&=1\cdot\det X(\vect x)\cdot1\\
&=-|\vect x|^2.
\end{align*}
Luego $|\Phi(U)\vect x|=|\vect x|$ para todo $\vect x\in\R^3$. Por lo tanto, $\Phi(U)\in O(3)$.

Falta ver que el determinante es $+1$. Recordemos que todo elemento de $\SU(2)$ puede escribirse en la forma
\[
U=\begin{pmatrix}\alpha&\beta\\-\beta^*&\alpha^*\end{pmatrix},
\qquad |\alpha|^2+|\beta|^2=1.
\]
Por tanto, como variedad real, $\SU(2)$ se identifica con la esfera $S^3$, y en particular es conexo. Como $\Phi$ es continua, la función
\[
U\longmapsto\det\Phi(U)
\]
toma valores en el conjunto discreto $\{+1,-1\}$. Por continuidad, no puede cambiar de signo dentro de $\SU(2)$. Como $\Phi(I_2)=I_3$, se tiene $\det\Phi(I_2)=1$. Así,
\[
\det\Phi(U)=1,\qquad \forall U\in\SU(2).
\]
Concluimos que $\Phi(U)\in\SO(3)$.

\subsection*{$\Phi$ es homomorfismo}

Sean $U,V\in\SU(2)$. Entonces
\begin{align*}
X\bigl(\Phi(UV)\vect x\bigr)
&=(UV)X(\vect x)(UV)^\dagger\\
&=UVX(\vect x)V^\dagger U^\dagger\\
&=U\bigl(VX(\vect x)V^\dagger\bigr)U^\dagger\\
&=UX\bigl(\Phi(V)\vect x\bigr)U^\dagger\\
&=X\bigl(\Phi(U)\Phi(V)\vect x\bigr).
\end{align*}
Como la correspondencia $\vect x\mapsto X(\vect x)$ es inyectiva, se obtiene
\[
\Phi(UV)\vect x=\Phi(U)\Phi(V)\vect x,\qquad \forall\vect x\in\R^3.
\]
Luego $\Phi(UV)=\Phi(U)\Phi(V)$. Por lo tanto, $\Phi$ es un homomorfismo de grupos.

\subsection*{Cálculo del núcleo}

El núcleo de $\Phi$ está dado por
\[
\kerp\Phi=\set{U\in\SU(2):\Phi(U)=I_3}.
\]
Si $U\in\kerp\Phi$, entonces $X(\Phi(U)\vect x)=X(\vect x)$ para todo $\vect x\in\R^3$. Pero, por definición de $\Phi$,
\[
UX(\vect x)U^\dagger=X(\vect x),\qquad \forall\vect x\in\R^3.
\]
Multiplicando a la derecha por $U$, se obtiene
\[
UX(\vect x)=X(\vect x)U,\qquad \forall\vect x\in\R^3.
\]
En particular, $U$ conmuta con las tres matrices de Pauli: $U\sigma_i=\sigma_iU$, $i=1,2,3$.

Escribamos
\[
U=\begin{pmatrix}a&b\\c&d\end{pmatrix}.
\]
Primero, de $U\sigma_3=\sigma_3U$, se tiene
\[
\begin{pmatrix}a&-b\\c&-d\end{pmatrix}
=\begin{pmatrix}a&b\\-c&-d\end{pmatrix}.
\]
Por lo tanto, $b=0$, $c=0$, y así
\[
U=\begin{pmatrix}a&0\\0&d\end{pmatrix}.
\]
Ahora, de $U\sigma_1=\sigma_1U$, se obtiene
\[
\begin{pmatrix}0&a\\d&0\end{pmatrix}
=\begin{pmatrix}0&d\\a&0\end{pmatrix}.
\]
Luego $a=d$ y, por tanto, $U=aI_2$. Como $U\in\SU(2)$, se cumple $\det U=1$. Pero $\det(aI_2)=a^2$. Entonces $a^2=1$, y de aquí $a=\pm1$. Por lo tanto,
\[
\kerp\Phi=\set{I_2,-I_2}=\mathbb{Z}_2.
\]

\subsection*{Sobreyectividad de $\Phi$}

Mostraremos ahora que toda rotación de $\SO(3)$ proviene de algún elemento de $\SU(2)$. Sea $\vect n\in\R^3$ un vector unitario, $|\vect n|=1$, y sea $\theta\in\R$. Definimos
\[
N=\vect n\cdot\vect\sigma.
\]
Como $|\vect n|=1$, se tiene $N^2=I_2$. Consideremos ahora la matriz
\[
U_{\vect n,\theta}=\cos\left(\frac{\theta}{2}\right)I_2-i\sin\left(\frac{\theta}{2}\right)N.
\]
Esta matriz pertenece a $\SU(2)$. En efecto,
\begin{align*}
U_{\vect n,\theta}^\dagger U_{\vect n,\theta}
&=\left(\cos\frac{\theta}{2}I_2+i\sin\frac{\theta}{2}N\right)
  \left(\cos\frac{\theta}{2}I_2-i\sin\frac{\theta}{2}N\right)\\
&=\cos^2\frac{\theta}{2}I_2+\sin^2\frac{\theta}{2}N^2\\
&=I_2.
\end{align*}
Además, como $N$ tiene autovalores $+1$ y $-1$, los autovalores de $U_{\vect n,\theta}$ son $e^{-i\theta/2}$ y $e^{i\theta/2}$. Por lo tanto,
\[
\det U_{\vect n,\theta}=e^{-i\theta/2}e^{i\theta/2}=1,
\]
y así $U_{\vect n,\theta}\in\SU(2)$.

Calculemos ahora la acción inducida sobre $\R^3$. Escribamos
\[
c=\cos\left(\frac{\theta}{2}\right),\qquad s=\sin\left(\frac{\theta}{2}\right).
\]
Entonces $U_{\vect n,\theta}=cI_2-isN$ y $U_{\vect n,\theta}^\dagger=cI_2+isN$. Para $X(\vect x)=\vect x\cdot\vect\sigma$, se tiene
\begin{align*}
U_{\vect n,\theta}X(\vect x)U_{\vect n,\theta}^\dagger
&=(cI_2-isN)X(\vect x)(cI_2+isN)\\
&=c^2X(\vect x)+ics\bigl(X(\vect x)N-NX(\vect x)\bigr)+s^2NX(\vect x)N.
\end{align*}
Usando la identidad de Pauli, se obtiene
\[
X(\vect x)N-NX(\vect x)=2i(\vect x\times\vect n)\cdot\vect\sigma=-2i(\vect n\times\vect x)\cdot\vect\sigma,
\]
y también
\[
NX(\vect x)N=\bigl(2(\vect n\cdot\vect x)\vect n-\vect x\bigr)\cdot\vect\sigma.
\]
Luego
\[
U_{\vect n,\theta}X(\vect x)U_{\vect n,\theta}^\dagger
=\bigl[(c^2-s^2)\vect x+2cs(\vect n\times\vect x)+2s^2(\vect n\cdot\vect x)\vect n\bigr]\cdot\vect\sigma.
\]
Como $c^2-s^2=\cos\theta$, $2cs=\sin\theta$ y $2s^2=1-\cos\theta$, resulta
\[
U_{\vect n,\theta}X(\vect x)U_{\vect n,\theta}^\dagger
=X\bigl(\cos\theta\,\vect x+\sin\theta(\vect n\times\vect x)+(1-\cos\theta)(\vect n\cdot\vect x)\vect n\bigr).
\]
La expresión vectorial
\[
R_{\vect n,\theta}\vect x=\cos\theta\,\vect x+\sin\theta(\vect n\times\vect x)+(1-\cos\theta)(\vect n\cdot\vect x)\vect n
\]
es la fórmula de Rodrigues para una rotación en $\R^3$ de ángulo $\theta$ alrededor del eje $\vect n$. Por lo tanto,
\[
\Phi(U_{\vect n,\theta})=R_{\vect n,\theta}.
\]

Pero todo elemento de $\SO(3)$ es una rotación de este tipo. En efecto, si $R\in\SO(3)$, entonces $R$ es una matriz ortogonal real de dimensión tres y determinante uno. Sus autovalores tienen módulo uno y, al ser el polinomio característico de grado impar, existe al menos un autovalor real. La condición $\det R=1$ fuerza que ese autovalor real sea $1$. Por lo tanto, existe un vector unitario $\vect n$ tal que $R\vect n=\vect n$. El plano perpendicular a $\vect n$ queda invariante y la restricción de $R$ a dicho plano es una rotación plana de cierto ángulo $\theta$. Así, existe un eje unitario $\vect n$ y un ángulo $\theta$ tales que $R=R_{\vect n,\theta}$.

En consecuencia, para todo $R\in\SO(3)$ existe $U_{\vect n,\theta}\in\SU(2)$ tal que $\Phi(U_{\vect n,\theta})=R$. Así, $\im\Phi=\SO(3)$ y $\Phi$ es sobreyectiva.

\subsection*{Aplicación del teorema fundamental de los homomorfismos}

Hemos construido un homomorfismo sobreyectivo
\[
\Phi:\SU(2)\longrightarrow\SO(3),
\]
con núcleo
\[
\kerp\Phi=\set{I_2,-I_2}=\mathbb{Z}_2.
\]
Por el teorema fundamental de los homomorfismos,
\[
\SU(2)/\kerp\Phi\cong\im\Phi.
\]
Como $\im\Phi=\SO(3)$, se concluye que
\[
\SU(2)/\mathbb{Z}_2\cong\SO(3).
\]

\subsection*{Resultado}

El homomorfismo
\[
\Phi:\SU(2)\longrightarrow\SO(3)
\]
definido por
\[
X\bigl(\Phi(U)\vect x\bigr)=UX(\vect x)U^\dagger
\]
es sobreyectivo y tiene núcleo
\[
\kerp\Phi=\set{I_2,-I_2}\cong\mathbb{Z}_2.
\]
Por lo tanto,
\[
\SU(2)/\mathbb{Z}_2\cong\SO(3).
\]
Esto significa que $\SU(2)$ cubre dos veces a $\SO(3)$: los elementos $U$ y $-U$ de $\SU(2)$ representan la misma rotación en $\SO(3)$.

\end{document}
